% \iffalse meta-comment
% A LaTeX Class for Semantic Lectures Slides (originally developed for Michael Kohlhase)
% Copyright (c) 2019 Michael Kohlhase, all rights reserved
%               this file is released under the
%               Gnu Library Public Licences (LGPL)
%
% The original of this file is in the public repository at 
% http://github.com/sLaTeX/sTeX/
% \fi
% 
% \iffalse
%<cls|package>\NeedsTeXFormat{LaTeX2e}[1999/12/01]
%<cls>\ProvidesClass{mikoslides}[2020/10/17 v1.3 MiKo slides Class]
%<package>\ProvidesPackage{mikoslides}[2020/10/17 v1.3 MiKo slides Package]
%
%<*driver>
\documentclass{ltxdoc}
\usepackage[utf8]{inputenc}
\usepackage{url,array,float,xspace}
\usepackage[show]{ed}
\usepackage{ctangit}
\usepackage{graphicx,paralist}
\usepackage[hyperref=auto,style=alphabetic]{biblatex}
\addbibresource{kwarcpubs.bib}
\addbibresource{extpubs.bib}
\addbibresource{kwarccrossrefs.bib}
\addbibresource{extcrossrefs.bib}
\makeindex
\floatstyle{boxed}
\newfloat{exfig}{thp}{lop}
\floatname{exfig}{Example}
\def\githubissue#1{\cite{sTeX:github:on},
\hyperlink{https://github.com/sLaTeX/sTeX/issues/#1}{issue #1}}
\usepackage{hyperref}
\begin{document}
\RecordChanges
\DocInput{mikoslides.dtx}
\end{document}
%</driver>
% \fi
% 
% \iffalse\CheckSum{469}\fi
% 
% \changes{v0.1}{2005/12/06}{Initial Version}
% \changes{v0.2}{2006/01/11}{course notes back on seminar}
% \changes{v0.3}{2007/05/16}{changing to Jacobs logo}
% \changes{v0.4}{2007/10/16}{re-basing the whole thing on beamer}
% \changes{v0.4}{2008/09/06}{moving line-end-comment to \texttt{omdoc.dtx}}
% \changes{v0.5}{2009/06/17}{eliminating mytwocolumns, this is better done by \texttt{beamer.cls}}
% \changes{v0.9}{2010/06/18}{this is almost done}
% \changes{v0.9}{2012/09/17}{basic options handling for the \texttt{frame} environment in
% notes mode}
% \changes{v0.9}{2013/08/23}{numbered sectocframes}
% \changes{v1.0}{2014/01/07}{adding \texttt{\textbackslash frameimage}}
% \changes{v1.1}{2015/10/25}{Removing the old title macros (use the regular ones instead)}
% \changes{v1.1}{2015/10/25}{reinterpreting omgroup}
% \changes{v1.1}{2015/11/04}{moving MathHub support out to separate package}
% \changes{v1.2}{2018/12/03}{changed to keyval class/package options, allowed arbitrary classes}
% \changes{v1.3}{2020/10/17}{reusing the sectioning counters of beamer}
% 
% \GetFileInfo{mikoslides.cls}
% \MakeShortVerb{\|}
%
% \def\twin#1#2{\index{#1!#2}\index{#2!#1}}
% \def\twintoo#1#2{{#1 #2}\twin{#1}{#2}}
% \def\atwin#1#2#3{\index{#1!#2!#3}\index{#3!#2 (#1)}}
% \def\atwintoo#1#2#3{{#1 #2 #3}\atwin{#1}{#2}{#3}}
%
% \def\scsys#1{{{\sc #1}}\index{#1@{\sc #1}}}
% \def\stex{\hbox{\raisebox{-.5ex}S\kern-.5ex\TeX}}
% \def\sTeX{\stex}
% \def\cnxml{\scshape{CNXml}}
% \def\connexions{\scshape{Connexions}}
% \def\element#1{{\ttfamily{#1}}}
% \def\snippet#1{{\ttfamily{#1}}}
% \def\cnxlatex{CNX\LaTeX\xspace}
% \def\mathml{{\scshape{MathML}}\xspace}
% \def\omdoc{OMDoc\xspace}
% \def\activemath{{\scshape{ActiveMath}}\xspace}
% \def\textwarning{\includegraphics[width=1.2em]{dangerous-bend}\xspace}
% 
% \title{Slides and Course Notes\thanks{Version {\fileversion}
% (last revised {\filedate})}}
%    \author{Michael Kohlhase\\
%            FAU Erlangen-N\"urnberg\\
%            \url{http://kwarc.info/kohlhase}}
% \maketitle
%
% \begin{abstract}
%   We present a document class from which we can generate both course slides and course
%   notes in a transparent way.
% \end{abstract}
%
% \tableofcontents\newpage
%
%\section{Introduction}
%
% The |mikoslides| document class is derived from |beamer.cls|~\cite{beamerclass:on}, it
% adds a ``notes version'' for course notes derived from the |omdoc|
% class~\cite{Kohlhase:smomdl} that is more suited to printing than the one supplied by
% |beamer.cls|.
%
% 
%\section{The User Interface}\label{sec:user}
%
% The |mikoslides| class takes the notion of a slide frame from Till Tantau's excellent
% |beamer| class and adapts its notion of frames for use in the \sTeX and \omdoc. To
% support semantic course notes, it extends the notion of mixing frames and explanatory
% text, but rather than treating the frames as images (or integrating their contents into
% the flowing text), the |mikoslides| package displays the slides as such in the course
% notes to give students a visual anchor into the slide presentation in the course (and to
% distinguish the different writing styles in slides and course notes).
% 
% In practice we want to generate two documents from the same source: the slides for
% presentation in the lecture and the course notes as a narrative document for home
% study. To achieve this, the |mikoslides| class has two modes: \emph{slides mode} and
% \emph{notes mode} which are determined by the package option. 
%
% \subsection{Package Options}\label{sec:user:options}
% 
% The |mikoslides| class takes a variety of class options:\ednote{leaving out noproblems
% for the moment until we decide what to do with it.}
% \begin{itemize}
% \item The options \DescribeMacro{slides}|slides|\DescribeMacro and {notes}|notes| switch
%   between slides mode and notes mode (see Section~\ref{sec:user:notesslides}).
% \item If the option \DescribeMacro{sectocframes}|sectocframes| is given, then for the
%   |omgroup|s, special frames with the |omgroup| title (and number) are generated.
% \item \DescribeMacro{showmeta}|showmeta|. If this is set, then the metadata keys are
%   shown (see~\ctancite{Kohlhase:metakeys} for details and customization options).
% \item If the option \DescribeMacro{frameimages}|frameimages| is set, then slide mode
%   also shows the |\frameimage|-generated frames. If also the
%   \DescribeMacro{fiboxed}|fiboxed| option is given, the slides are surrounded by a box.
% \item \DescribeMacro{topsect}|topsect=|\meta{sect} can be used to specify the
%   top-level sectioning level; the default for \meta{sect} is |section|.
% \item The \DescribeMacro{mh}|mh| option turns on MathHub support; see
%   \ctancite{Kohlhase:mss}.
% \end{itemize}
% 
% \subsection{Notes and Slides}\label{sec:user:notesslides}
% 
% Slides are represented with the \DescribeEnv{frame}|frame| just like in the |beamer|
% class, see~\cite{Tantau:ugbc} for details. The |mikoslides| class adds the
% \DescribeEnv{note}|note| environment for encapsulating the course note
% fragments.\footnote{MK: it would be very nice, if we did not need this environment, and
% this should be possible in principle, but not without intensive LaTeX trickery. Hints to
% the author are welcome.} 
%
% \textwarning Note that it is essential to start and end the |notes| environment at the
% start of the line -- in particular, there may not be leading blanks -- else {\LaTeX}
% becomes confused and throws error messages that are difficult to decipher.
%
% \begin{exfig}
% \begin{verbatim}
% \ifnotes\maketitle\else
% \frame[noframenumbering]\maketitle\fi
% 
% \begin{note}
%   We start this course with ...
% \end{note}
%
% \begin{frame}
%   \frametitle{The first slide}
%   ...
% \end{frame}
% \begin{note}
%   ... and more explanatory text
% \end{note}
%
% \begin{frame}
%   \frametitle{The second slide}
%   ...
% \end{frame}
% ...
% \end{verbatim}
% \caption{A typical Course Notes File}\label{fig:notesfile}
% \end{exfig}
% 
% By interleaving the |frame| and |note| environments, we can build course notes as shown
% in Figure~\ref{fig:notesfile}.
%
% Note the use of the \DescribeMacro{\ifnotes}|\ifnotes| conditional, which allows
% different treatment between |notes| and |slides| mode -- manually setting |\notestrue|
% or |\notesfalse| is strongly discouraged however.
% 
% \textwarning: We need to give the title frame the |noframenumbering| option so that the
% frame numbering is kept in sync between the slides and the course notes.
%
% \textwarning: The |beamer| class recommends not to use the |allowframebreaks| option on
% frames (even though it is very convenient). This holds even more in the |mikoslides|
% case: At least in conjunction with |\newpage|, frame numbering behaves funnily (we have
% tried to fix this, but who knows). 
% 
% Sometimes, we want to integrate slides as images after all -- e.g. because we already
% have a PowerPoint presentation, to which we want to add \sTeX notes. In this case we can
% use \DescribeMacro{\frameimage}|\frameimage[|\meta{opt}|]{|\meta{path}|}|, where
% \meta{opt} are the options of |\includegraphics| from the |graphicx|
% package~\cite{CarRah:tpp99} and \meta{path} is the file path (extension can be left off
% like in |\includegraphics|). We have added the |label| key that allows to give a frame
% label that can be referenced like a regular |beamer| frame.\ednote{MK: the hyperref link
% does not seem to work yet. I wonder why but do not have the time to fix it.}
% 
% If we want to transclude a the contents of a file as a note, we can use the
% \DescribeMacro{\inputref*}|\inputref*| macro: |\inputref*{foo}| is equivalent to 
%\begin{verbatim}
% \begin{note}
%   \inputref{foo}
% \end{note}
% \end{verbatim}
% 
% There are some environments that tend to occur at the top-level of |note|
% environments. We make convenience versions of these: e.g. the
% \DescribeEnv{nomtext}|nomtext| environment is just an |omtext| inside a |note|
% environment (but looks nicer in the source, since it avoids one level of source
% indenting). Similarly, we have the \DescribeEnv{nomgroup}|nomgroup| environment.
%  
% \subsection{Header and Footer Lines}\label{sec:user:headfootlines}
%
% \subsection{Colors and Highlighting}\label{sec:user:highlighting}
% The \DescribeMacro{\textwarning}|\textwarning| macro generates a warning
% sign: \textwarning
%
% \subsection{Front Matter, Titles, etc}\label{sec:user:matter}
%
% \subsection{Miscellaneous}\label{sec:user:misc}
%
% \section{Limitations}\label{sec:limitations}
% 
% In this section we document known limitations. If you want to help alleviate them,
% please feel free to contact the package author. Some of them are currently discussed in
% the \sTeX GitHub repository~\cite{sTeX:github:on}. 
% \begin{enumerate}
% \item when option |book| which uses |\pagestyle{headings}| is given and semantic macros
%   are given in the |omgroup| titles, then they sometimes are not defined by the time the
%   heading is formatted. Need to look into how the headings are made. This is a problem
%   of the underlying |omdoc| package.
% \end{enumerate}
% 
% \StopEventually{\newpage\PrintIndex\newpage\PrintChanges\printbibliography}
% 
%\section{The Implementation}\label{sec:impl}
%
%\subsection{Class and Package Options}\label{sec:impl:init}
%
% We define some Package Options and switches for the |mikoslides| class and activate them
% by passing them on to |beamer.cls| and |omdoc.cls| and the |mikoslides| package. We pass
% the |nontheorem| option to the |statements| package when we are not in notes mode, since
% the |beamer| package has its own (overlay-aware) theorem environments. 
%
%    \begin{macrocode}
%<*cls>
\RequirePackage{stex-base}
\RequirePackage{kvoptions}
\RequirePackage{etoolbox}
\SetupKeyvalOptions{family=mks@cls,prefix=mks@cls@}
\DeclareStringOption[article]{class}
\AddToKeyvalOption*{class}{\PassOptionsToClass{class=\mks@cls@class}{omdoc}
  \ifdefstring{\mks@cls@class}{book}{\PassOptionsToPackage{defaulttopsect=part}{mikoslides}}{}
  \ifdefstring{\mks@cls@class}{report}{\PassOptionsToPackage{defaulttopsect=part}{mikoslides}}}{}
\DeclareBoolOption{notes}
\DeclareComplementaryOption{slides}{notes}
\DeclareDefaultOption{%
  \PassOptionsToClass{\CurrentOption}{omdoc}
  \PassOptionsToClass{\CurrentOption}{beamer}
  \PassOptionsToPackage{\CurrentOption}{mikoslides}}
\ProcessKeyvalOptions{mks@cls}
%</cls>
%    \end{macrocode}
% now we do the same for the |mikoslides| package. 
%    \begin{macrocode}
%<*package>
\RequirePackage{stex-base}
\RequirePackage{kvoptions}
\SetupKeyvalOptions{family=mks@sty,prefix=mks@sty@}
\DeclareStringOption{topsect}
\DeclareStringOption{defaulttopsect}
\DeclareBoolOption{mh}
\AddToKeyvalOption*{mh}{
  \PassOptionsToPackage{mh}{stex}
  \PassOptionsToPackage{mh}{smglom}
  \PassOptionsToPackage{mh}{tikzinput}}
\newif\ifnotes\notestrue
\DeclareBoolOption{notes}
\AddToKeyvalOption*{notes}{\notestrue\PassOptionsToPackage{notes}{statements}}
\DeclareComplementaryOption{slides}{notes}
\AddToKeyvalOption*{slides}{\notesfalse\PassOptionsToPackage{nontheorem}{statements}}
\DeclareBoolOption{sectocframes}
\DeclareBoolOption{frameimages}
\DeclareBoolOption{fiboxed}
\DeclareBoolOption{noproblems}
\DeclareDefaultOption{%
  \PassOptionsToPackage{\CurrentOption}{stex}
  \PassOptionsToPackage{\CurrentOption}{smglom}
  \PassOptionsToPackage{\CurrentOption}{tikzinput}}
\ProcessKeyvalOptions{mks@sty}
%    \end{macrocode}
% we give ourselves a macro |\@@topsect| that needs only be evaluated once, so that the
% |\ifdefstring| conditionals work below.
%    \begin{macrocode}
\ifx\mks@sty@topsect\@empty\edef\@@topsect{\mks@sty@defaulttopsect}
\else\edef\@@topsect{\mks@sty@topsect}\fi
%</package>
%    \end{macrocode}
%
% Depending on the options, we either load the |article|-based |omdoc| or the |beamer|
% class (and set some counters).
%    \begin{macrocode}
%<*cls>
\ifmks@cls@notes
  \LoadClass{omdoc}
\else
  \LoadClass[10pt,notheorems,xcolor={dvipsnames,svgnames}]{beamer}
  \newcounter{Item}
  \newcounter{paragraph}
  \newcounter{subparagraph}
  \newcounter{Hfootnote}
\fi
%    \end{macrocode}
% now it only remains to load the |mikoslides| package that does all the rest. 
%    \begin{macrocode}
\RequirePackage{mikoslides}
%</cls>
%    \end{macrocode}
% 
% In |notes| mode, we also have to make the |beamer|-specific things available to
% |article| via the |beamerarticle| package. We use options to avoid loading theorem-like
% environments, since we want to use our own from the $\sTeX$ packages.  The first batch
% of packages we want are loaded on |mikoslides.sty|. These are the general ones, we will
% load the \sTeX-specific ones after we have done some work (e.g. defined the counters
% |m*|). Only the |stex-logo| package is already needed now for the default theme.
%
%    \begin{macrocode}
%<*package>
\ifmks@sty@notes
\RequirePackage{a4wide}
\RequirePackage{marginnote}
\RequirePackage[dvipsnames,svgnames]{xcolor}
\if@latexml\else\RequirePackage{mdframed}\fi
\RequirePackage[noxcolor,noamsthm]{beamerarticle}
\fi
\RequirePackage{etoolbox}
\RequirePackage{amssymb}
\RequirePackage{amsmath}
\RequirePackage{comment}
\RequirePackage{textcomp}
\RequirePackage{url}
\RequirePackage{graphicx}
\RequirePackage{stex-logo}
\RequirePackage{pgf}
\ifmks@sty@notes
\RequirePackage[bookmarks,bookmarksopen,bookmarksnumbered,breaklinks,hidelinks]{hyperref}
\fi
%    \end{macrocode}
%
% finally, we require the |metakeys| package from \sTeX, so that we can use the
% |\addmetakey| mechanism.
%
%    \begin{macrocode}
\RequirePackage{metakeys}
%    \end{macrocode}
% 
% \subsection{Notes and Slides}\label{sec:impl:noteslides}
%
% For the lecture notes cases, we also provide the |\usetheme| macro that would otherwise
% come from the the |beamer| class. While the latter loads
% |beamertheme|\meta{theme}{.sty}, the notes version loads
% |beamernotestheme|\meta{theme}|.sty|.\ednote{MK: This is not ideal, but I am not sure
% that I want to be able to provide the full theme functionality there.}
%    \begin{macrocode}
\ifmks@sty@notes
\renewcommand\usetheme[2][]{\usepackage[#1]{beamernotestheme#2}}
\fi
%    \end{macrocode}
% We define the sizes of slides in the notes. Somehow, we cannot get by with the same
% here. 
%
%    \begin{macrocode}
\newcounter{slide}
\newlength{\slidewidth}\setlength{\slidewidth}{13.5cm}
\newlength{\slideheight}\setlength{\slideheight}{9cm}
%    \end{macrocode}
% 
% \begin{environment}{note}
% The |note| environment is used to leave out text in the |slides| mode. It does not have
% a counterpart in OMDoc. So for course notes, we define the |note| environment to be a
% no-operation otherwise we declare the |note| environment as a comment via the |comment|
% package.
%    \begin{macrocode}
\ifmks@sty@notes%
  \renewenvironment{note}{\ignorespaces}{}%
\else%
  \excludecomment{note}%
\fi%
%    \end{macrocode}
% \end{environment}
% 
% We first set up the slide boxes in |article| mode. We set up sizes and provide a
% box register for the frames and a counter for the slides.
% 
%    \begin{macrocode}
\ifmks@sty@notes
  \newlength{\slideframewidth}
  \setlength{\slideframewidth}{1.5pt}
%    \end{macrocode}
% 
% \begin{environment}{frame}
%   We first define the keys. 
%    \begin{macrocode}
  \addmetakey{frame}{label}
  \addmetakey[yes]{frame}{allowframebreaks}
  \addmetakey{frame}{allowdisplaybreaks}
  \addmetakey[yes]{frame}{fragile}
  \addmetakey[yes]{frame}{shrink}
  \addmetakey[yes]{frame}{squeeze}
  \addmetakey[yes]{frame}{t}
%    \end{macrocode}
% We define the environment, read them, and construct the slide number and label.
%    \begin{macrocode}
  \renewenvironment{frame}[1][]{%
    \metasetkeys{frame}{#1}%
    \sffamily%
    \stepcounter{slide}%
    \def\@currentlabel{\theslide}%
    \ifx\frame@label\@empty%
    \else%
      \label{\frame@label}%
    \fi%
%    \end{macrocode}
%   We redefine the |itemize| environment so that it looks more like the one in |beamer|. 
%    \begin{macrocode}
    \def\itemize@level{outer}%
    \def\itemize@outer{outer}%
    \def\itemize@inner{inner}%
    \renewcommand\newpage{\addtocounter{framenumber}{1}}%
    \renewcommand\metakeys@show@keys[2]{\marginnote{{\scriptsize ##2}}}%
    \renewenvironment{itemize}{%
      \ifx\itemize@level\itemize@outer%
        \def\itemize@label{$\rhd$}%
      \fi%
      \ifx\itemize@level\itemize@inner%
        \def\itemize@label{$\scriptstyle\rhd$}%
      \fi%
      \begin{list}%
      {\itemize@label}%
      {\setlength{\labelsep}{.3em}%
       \setlength{\labelwidth}{.5em}%
       \setlength{\leftmargin}{1.5em}%
      }%
      \edef\itemize@level{\itemize@inner}%
    }{%
      \end{list}%
    }%
%    \end{macrocode}
% We create the box with the |mdframed| environment from the equinymous package.
%    \begin{macrocode}
    \begin{mdframed}[linewidth=\slideframewidth,skipabove=1ex,skipbelow=1ex,userdefinedwidth=\slidewidth,align=center]\sf%
  }{%
    \medskip\miko@slidelabel\end{mdframed}%
  }%
%    \end{macrocode}
% \end{environment}
% 
% Now, we need to redefine the frametitle (we are still in course notes mode). 
% \begin{macro}{\frametitle}
%    \begin{macrocode}
  \renewcommand{\frametitle}[1]{{\Large\bf\sf\color{blue}{#1}}\medskip}%
\fi %ifmks@sty@notes
%    \end{macrocode}
% \end{macro}
% 
% \begin{macro}{\frameimage}
%   We have to make sure that the width is overwritten, for that we check the
%   |\Gin@ewidth| macro from the |graphicx| package. We also add the |label| key. 
%    \begin{macrocode}
\define@key{Gin}{label}{\def\@currentlabel{\arabic{slide}}\label{#1}}
\newrobustcmd\frameimage[2][]{%
  \stepcounter{slide}%
  \ifmks@sty@frameimages%
    \def\Gin@ewidth{}\setkeys{Gin}{#1}%
    \ifmks@sty@notes\else\vfill\fi%
    \begin{center}
      \ifmks@sty@fiboxed%
        \fbox{\ifx\Gin@ewidth\@empty\includegraphics[width=\slidewidth,#1]{#2}\else\mygraphics[#1]{#2}\fi}%
      \else
        \ifx\Gin@ewidth\@empty\includegraphics[width=\slidewidth,#1]{#2}\else\mygraphics[#1]{#2}\fi%
      \fi% ifmks@sty@fiboxed
     \end{center}
    \par\strut\hfill{\footnotesize Slide \arabic{slide}}%
    \ifmks@sty@notes\else\vfill\fi%
  \fi} % ifmks@sty@frameimages
%    \end{macrocode}
% \end{macro}
% 
% \begin{macro}{\pause}
%   \ednote{MK: fake it in notes mode for now}
%    \begin{macrocode}
\ifmks@sty@notes\newcommand\pause{}\fi
%    \end{macrocode}
% \end{macro}
% 
% \begin{environment}{nomtext}
%    \begin{macrocode}
\ifmks@sty@notes\newenvironment{nomtext}[1][]{\begin{omtext}[#1]}{\end{omtext}}%
\else\excludecomment{nomtext}\fi%
%    \end{macrocode}
% \end{environment}
%
% \begin{environment}{nomgroup}
%    \begin{macrocode}
\ifmks@sty@notes\newenvironment{nomgroup}[2][]{\begin{omgroup}[#1]{#2}}{\end{omgroup}}%
\else\excludecomment{nomgroup}\fi%
%    \end{macrocode}
% \end{environment}
%
% \subsection{Header and Footer Lines}\label{sec:impl:headfootlines}
%
% Now, we set up the infrastructure for the footer line of the slides, we use boxes for
% the logos, so that they are only loaded once, that considerably speeds up processing.
% 
% \begin{macro}{\setslidelogo}
% The default logo is the \sTeX logo. Customization can be done by |\setslidelogo{|\meta{logo name}|}|.
%    \begin{macrocode}
\newlength{\slidelogoheight}
\ifmks@sty@notes%
  \setlength{\slidelogoheight}{.4cm}%
\else%
  \setlength{\slidelogoheight}{1cm}%
\fi%
\newsavebox{\slidelogo}%
\sbox{\slidelogo}{\sTeX}%
\newrobustcmd{\setslidelogo}[1]{%
  \sbox{\slidelogo}{\includegraphics[height=\slidelogoheight]{#1}}%
}%
%    \end{macrocode}
% \end{macro}
%
% \begin{macro}{\setsource}
% |\source| stores the writer's name. By default it is {\it Michael Kohlhase} since he is the main user and designer of this package. |\setsource{|\meta{name}|}| can change the writer's name.
%    \begin{macrocode}
\def\source{Michael Kohlhase}% customize locally
\newrobustcmd{\setsource}[1]{\def\source{#1}}%
%    \end{macrocode}
% \end{macro}
%
% \begin{macro}{\setlicensing}
%   Now, we set up the copyright and licensing. By default we use the Creative Commons
%   Attribuition-ShareAlike license to strengthen the public domain. If package |hyperref|
%   is loaded, then we can attach a hyperlink to the license
%   logo. |\setlicensing[|\meta{url}|]{|\meta{logo name}|}| is used for customization,
%   where ||\meta{url}|| is optional.
%    \begin{macrocode}
\def\copyrightnotice{\footnotesize\copyright:\hspace{.3ex}{\source}}%
\newsavebox{\cclogo}%
\sbox{\cclogo}{\includegraphics[height=\slidelogoheight]{cc_somerights}}%
\newif\ifcchref\cchreffalse%
\AtBeginDocument{%
  \@ifpackageloaded{hyperref}{\cchreftrue}{\cchreffalse}
}%
\def\licensing{%
  \ifcchref%
    \href{http://creativecommons.org/licenses/by-sa/2.5/}{\usebox{\cclogo}}%
  \else%
    {\usebox{\cclogo}}%
  \fi%
}%
\newrobustcmd{\setlicensing}[2][]{%
  \def\@url{#1}%
  \sbox{\cclogo}{\includegraphics[height=\slidelogoheight]{#2}}%
  \ifx\@url\@empty%
    \def\licensing{{\usebox{\cclogo}}}%
  \else%
    \def\licensing{%
      \ifcchref%
      \href{#1}{\usebox{\cclogo}}%
      \else%
      {\usebox{\cclogo}}%
      \fi%
    }%
  \fi%
}%
%    \end{macrocode}
% \end{macro} 
%
% \begin{macro}{\slidelabel}
% Now, we set up the slide label for the |article| mode.\ednote{see that we can use the themes for the slides some day. This is all fake.}
%    \begin{macrocode}
\newrobustcmd\miko@slidelabel{%
  \vbox to \slidelogoheight{%
    \vss\hbox to \slidewidth%
    {\licensing\hfill\copyrightnotice\hfill\arabic{slide}\hfill\usebox{\slidelogo}}%
  }%
}%
%    \end{macrocode}
% \end{macro}
% 
% \subsection{Colors and Highlighting}\label{sec:impl:highlighting}
%
% We first specify sans serif fonts as the default. 
%
%    \begin{macrocode}
\sffamily
%    \end{macrocode}
%
% Now, we set up an infrastructure for highlighting phrases in slides. Note that we use
% content-oriented macros for highlighting rather than directly using color markup. 
% The first thing to to is to adapt the green so that it is dark enough for most beamers
%    \begin{macrocode}
\AtBeginDocument{%
\definecolor{green}{rgb}{0,.5,0}%
\definecolor{purple}{cmyk}{.3,1,0,.17}%
}%
%    \end{macrocode}
%
% We customize the |\defemph|, |\notemph|, and |\stDMemph| macros with colors for the use
% in the |statements| package. Furthermore we customize the |\@@lec| macro for the
% appearance of line end comments in |\lec|.
%
%    \begin{macrocode}
% \def\STpresent#1{\textcolor{blue}{#1}}
\def\defemph#1{{\textcolor{magenta}{#1}}}
\def\termemph#1{{\textcolor{cyan}{#1}}}
\def\notemph#1{{\textcolor{magenta}{#1}}}
\def\stDMemph#1{{\textcolor{blue}{#1}}}
\def\@@lec#1{(\textcolor{green}{#1})}
%    \end{macrocode}
%
% I like to use the dangerous bend symbol for warnings, so we provide it here.
% \begin{macro}{\textwarning}
%   as the macro can be used quite often we put it into a box register, so that it is only
%   loaded once. 
%    \begin{macrocode}
\pgfdeclareimage[width=.8em]{miko@small@dbend}{dangerous-bend}
\def\smalltextwarning{%
  \pgfuseimage{miko@small@dbend}%
  \xspace%
}%
\pgfdeclareimage[width=1.2em]{miko@dbend}{dangerous-bend}
\newrobustcmd\textwarning{%
  \raisebox{-.05cm}{\pgfuseimage{miko@dbend}}%
  \xspace%
}%
\pgfdeclareimage[width=2.5em]{miko@big@dbend}{dangerous-bend}%
\newrobustcmd\bigtextwarning{%
  \raisebox{-.05cm}{\pgfuseimage{miko@big@dbend}}%
  \xspace%
}%
%    \end{macrocode}
% \end{macro}
% 
%    \begin{macrocode}
\newrobustcmd\putgraphicsat[3]{%
  \begin{picture}(0,0)\put(#1){\includegraphics[#2]{#3}}\end{picture}%
}%
\newrobustcmd\putat[2]{%
  \begin{picture}(0,0)\put(#1){#2}\end{picture}%
}%
%    \end{macrocode}
%
% \subsection{Sectioning}
%
% If the |sectocframes| option is set, then we make section frames. We first define
% counters for |part| and |chapter|, which |beamer.cls| does not have and we make the
% |section| counter which it does dependent on |chapter|. 
%    \begin{macrocode}
\ifmks@sty@sectocframes%
\ifdefstring\@@topsect{part}{%
  \newcounter{chapter}\counterwithin*{section}{chapter}}
{\ifdefstring\@@topsect{chapter}{\newcounter{chapter}\counterwithin*{section}{chapter}}{}}
\fi% ifsectocframes
%    \end{macrocode}
%
% Now that we have defined the counters, we can load the \sTeX-specific packages (in
% particular |statements| that needs these counters). 
%
%    \begin{macrocode}
\RequirePackage{stex}
\RequirePackage{smglom}
\RequirePackage{tikzinput}
\ifmks@sty@mh\RequirePackage{mikoslides-mh}\fi
%    \end{macrocode}
% 
% Finally, we set the \DescribeMacro{\section@level}|\section@level| macro that governs
% sectioning. 
%
%    \begin{macrocode}
\section@level=2
\ifdefstring{\@@topsect}{part}
{\section@level=0%
  \def\thesection{\arabic{chapter}.\arabic{section}}%
  \def\part@prefix{\arabic{chapter}.}}{}
\ifdefstring{\@@topsect}{chapter}
{\section@level=1%
  \def\thesection{\arabic{chapter}.\arabic{section}}%
  \def\part@prefix{\arabic{chapter}.}}{}
\ifmks@sty@notes\else% only in slides
\renewenvironment{omgroup}[2][]{%
  \metasetkeys{omgroup}{#1}\sref@target%
  \advance\section@level by 1%
  \advance\omgroup@level by 1%
  \ifmks@sty@sectocframes%
  \stepcounter{slide}
  \begin{frame}[noframenumbering]%
  \vfill\Large\centering%
  \red{%
    \ifcase\section@level\or
    \stepcounter{part}
    \def\@@label{\omdoc@part@kw~\Roman{part}}
    \def\currentsectionlevel{\omdoc@part@kw}
    \or%
    \stepcounter{chapter}
    \def\@@label{\omdoc@chapter@kw~\arabic{chapter}}
    \def\currentsectionlevel{\omdoc@chapter@kw}
    \or
    \stepcounter{section}
    \def\@@label{\part@prefix\arabic{section}}
    \def\currentsectionlevel{\omdoc@section@kw}
    \or
    \stepcounter{subsection}
    \def\@@label{\part@prefix\arabic{section}.\arabic{subsection}}
    \def\currentsectionlevel{\omdoc@subsection@kw}
    \or
    \stepcounter{subsubsection}
    \def\@@label{\part@prefix\arabic{section}.\arabic{subsection}.\arabic{subsubsection}}
    \def\currentsectionlevel{\omdoc@subsubsection@kw}
    \or
    \stepcounter{mparagraph}
    \def\@@label{\part@prefix\arabic{section}.\arabic{msubsection}.\arabic{subsubsection}.\arabic{mparagraph}}
    \def\currentsectionlevel{\omdoc@paragraph@kw}
    \fi% end ifcase
    \@@label\sref@label@id\@@label
    \quad #2%
  }%
  \vfill%
  \end{frame}%
  \fi %ifmks@sty@sectocframes
}
{\advance\section@level by -1}%
\fi% ifmks@sty@notes
%    \end{macrocode}
%
% \begin{macro}{\inputref@*skip}
% We customize the hooks for in |\inputref|. 
%    \begin{macrocode}
\def\inputref@preskip{\smallskip}
\def\inputref@postskip{\medskip}
%    \end{macrocode}
% \end{macro}
%
% \begin{macro}{\inputref*}
%    \begin{macrocode}
\let\orig@inputref\inputref
\def\inputref{\@ifstar\ninputref\orig@inputref}
\newcommand\ninputref[2][]{\ifmks@sty@notes\orig@inputref[#1]{#2}\fi}
%    \end{macrocode}
% \end{macro}
% 
% \subsection{Miscellaneous}
%
% We set up a |beamer| template for theorems like ams style, but without a block
% environment.  
%    \begin{macrocode}
\def\inserttheorembodyfont{\normalfont}
\ifmks@sty@notes\else% only in slides
\defbeamertemplate{theorem begin}{miko}
{\inserttheoremheadfont\inserttheoremname\inserttheoremnumber
  \ifx\inserttheoremaddition\@empty\else\ (\inserttheoremaddition)\fi%
  \inserttheorempunctuation\inserttheorembodyfont\xspace}
\defbeamertemplate{theorem end}{miko}{}
%    \end{macrocode}
% and we set it as the default one. 
%    \begin{macrocode}
\setbeamertemplate{theorems}[miko]
%    \end{macrocode}
% The following fixes an error I do not understand, this has something to do with
% beamer compatibility, which has similar definitions but only up to 1. 
%    \begin{macrocode}
\expandafter\def\csname Parent2\endcsname{}
\fi% ifmks@sty@notes
\ifmks@sty@notes%
  \renewenvironment{columns}[1][]{%
    \par\noindent%
    \begin{minipage}%
    \slidewidth\centering\leavevmode%
  }{%
    \end{minipage}\par\noindent%
  }%
  \newsavebox\columnbox%
  \renewenvironment<>{column}[2][]{%
    \begin{lrbox}{\columnbox}\begin{minipage}{#2}%
  }{%
    \end{minipage}\end{lrbox}\usebox\columnbox%
  }%
\fi% ifmks@sty@notes
%    \end{macrocode}
%
%    \begin{macrocode}
\ifmks@sty@noproblems%
  \newenvironment{problems}{}{}%
\else%
  \excludecomment{problems}%
\fi%
%</package>
%    \end{macrocode}
% \Finale
\endinput
% \iffalse
% LocalWords:  mikoslides dtx beamer omdoc notheorems noamsthm beamerarticle sc Licences
% LocalWords:  graphicx slidelabel stex amssymb tikz url CPERL  amsmath filedate keyval
% LocalWords:  LoadClass RequirePackage DefRegister DefEnvironment omgroup rgb frameimage
% LocalWords:  afterDigestBegin setProperty LookupValue DefConstructor hyperref Tantau's
% LocalWords:  cmyk lec DefMacro titleslide ttitle RawTeX metadata etoolbox cls emph ugbc
% LocalWords:  noproblems linkcolor bookmarksopen citecolor urlcolor colorlinks Tantau di
% LocalWords:  breaklinks plainpages pdfpagelabels srcref iffalse texttt atwin exfig dt
% LocalWords:  mytwocolumns twintoo atwintoo scsys sc hbox raisebox cnxml impl notesfile
% LocalWords:  scshape ttfamily cnxlatex mathml activemath fileversion newpage tpp99 dd
% LocalWords:  maketitle tableofcontents ednote compactitem showmeta showmeta beamerclass
% LocalWords:  sectocframes sectocframes textwarning textwarning compactenum eq omgroups
% LocalWords:  includegraphics tracissue printbibliography textsf langle textsf rangle
% LocalWords:  langle ltxml metakeys newif ifnotes notesfalse ifsectocframes rangle putat
% LocalWords:  sectocframesfalse ifproblems problemstrue notestrue marginnote frontmatter
% LocalWords:  problemsfalse sectocframestrue mdframed noxcolor newcounter ifx equinymous
% LocalWords:  Hfootnote usetheme tikzinput usepgflibrary usetikzlibrary rhd ignorespaces
% LocalWords:  tikzmark textcomp newlength slidewidth setlength slidewidth miko setkeys
% LocalWords:  slideheight slideheight renewenvironment excludecomment itenize Gin@ewidth
% LocalWords:  slideframewidth slideframewidth surroundwithmdframed addmetakey Gen@ewidth
% LocalWords:  allowframebreaks allowdisplaybreaks metasetkeys stepcounter sbox mpart baz
% LocalWords:  currentlabel theslide renewcommand scriptsize scriptstyle hspace mchapter
% LocalWords:  medskip linewidth skipabove skipbelow frametitle newenvironment msection
% LocalWords:  slidelogoheight newsavebox slidelogo slidelogo vbox ifcase
% LocalWords:  Attribuition-ShareAlike copyrightnotice footnotesize cclogo vss minipage
% LocalWords:  cclogo somerights ifcchref cchreffalse ifpackageloaded usebox mycgraphics
% LocalWords:  cchreftrue usebox newcommand hfill hfill definecolor definecolor endinput
% LocalWords:  defemph notemph stDMemph STpresent textcolor textwarigrening ltx HorIacJuc
% LocalWords:  pgfdeclareimage dbend pgfuseimage xspace titleframe titlepage mycgraphics
% LocalWords:  titleframewith hline vspace ttitlejoint newbox boxwith boxwith msubsection
% LocalWords:  putgraphicsat beginomgroup ifnum vfill vfill noindent leavevmode cscpnrr11
% LocalWords:  columnbox lrbox needwrapper unlist omtext bgroup autoclose pgf readXToken
% LocalWords:  mygraphics doctex NeedsTeXFormat textbackslash userdefinedwidth foobar
%  LocalWords:  includegrahics smalltextwarnings msubsubsection mparagraph ldots
%  LocalWords:  frameimages frameimages ifframeimages frameimagesfalse frameimagestrue
%  LocalWords:  expandafter csname endcsname specializes customization flexiformal colors
%  LocalWords:  initialize Initialization defindex realized itemize newrobustcmd fiboxed
%  LocalWords:  specialize centering itemizations setsource setlicensing Kohlhase:smomdl
%  LocalWords:  fiboxed topsect topsect transclude inputref inputref inputref smglom
%  LocalWords:  nomtext nomtext nomgroup nomgroup nontheorem kvoptions mks@cls,prefix
%  LocalWords:  mks@cls ifdefstring mks@sty,prefix mks@sty 10pt,notheorems,xcolor xcolor
%  LocalWords:  dvipsnames,svgnames noxcolor,noamsthm black,citecolor black,urlcolor
%  LocalWords:  bookmarks,bookmarksopen,bookmarksnumbered,breaklinks cyan,filecolor
%  LocalWords:  cyan,colorlinks beamertheme beamernotestheme usepackage sffamily
%  LocalWords:  mycbgraphics mycbgraphics smalltextwarning bigtextwarning sref@target
%  LocalWords:  noframenumbering currentsectionlevel subsubsection paragraph
%  LocalWords:  ninputref inserttheorembodyfont defbeamertemplate inserttheoremheadfont
%  LocalWords:  inserttheoremname inserttheoremnumber inserttheoremaddition
%  LocalWords:  inserttheoremaddition inserttheorempunctuation setbeamertemplate
% \fi
% \endinput
% Local Variables:
% mode: doctex
% TeX-master: t
% End:
